\subsubsection*{Численный пример}

\begin{idea}[title={Откуда этот материал}]
Прогон метода Хука--Дживса на числах; на лекции метод был описан только
схемой.
\end{idea}

\begin{example}\label{ex:hooke-jeeves}
\[
  f(x_1, x_2) = (x_1 - 1)^2 + (x_2 - 2)^2 \ \longrightarrow\ \min ,
  \qquad x \in \R^2 .
\]
Минимум очевиден — $x_* = (1, 2)\T$, но метод до него доберётся, не вычисляя ни
одной производной. Начальные данные:
\[
  x^{(0)} = \begin{pmatrix} 0 \\ 0 \end{pmatrix} , \qquad
  h_1 = h_2 = 1 , \qquad \beta = \tfrac12 , \qquad \eps = 0{,}1 ,
\]
\[
  f\big(x^{(0)}\big) = (0-1)^2 + (0-2)^2 = 5 .
\]

\medskip
\noindent\textbf{Итерация 1. Исследующий поиск из $x^{(0)} = (0,0)\T$.}

Координата $x_1$, пробный шаг вперёд:
\[
  (0,0)\T + 1\cdot e_1 = (1,0)\T , \qquad f(1,0) = 0 + 4 = 4 < 5 ,
\]
шаг принимается, текущая точка $(1,0)\T$.

Координата $x_2$, пробный шаг вперёд:
\[
  (1,0)\T + 1\cdot e_2 = (1,1)\T , \qquad f(1,1) = 0 + 1 = 1 < 4 ,
\]
шаг принимается. Координаты кончились:
\[
  x^{(1)} = (1,1)\T , \qquad f\big(x^{(1)}\big) = 1 < 5 ,
\]
исследующий поиск удачен.

\medskip
\noindent\textbf{Поиск по образцу.}
\[
  x_{\text{обр}} = 2x^{(1)} - x^{(0)} = 2\begin{pmatrix} 1 \\ 1 \end{pmatrix}
    - \begin{pmatrix} 0 \\ 0 \end{pmatrix}
    = \begin{pmatrix} 2 \\ 2 \end{pmatrix} ,
  \qquad
  f(2,2) = 1 + 0 = 1 .
\]
Убывания нет: $f(x_{\text{обр}}) = 1 \not< f\big(x^{(1)}\big) = 1$ — точка по
образцу отвергается, новой базовой остаётся $x^{(1)} = (1,1)\T$.

\medskip
\noindent\textbf{Итерация 2. Исследующий поиск из $x^{(1)} = (1,1)\T$.}

Координата $x_1$, вперёд:
\[
  (2,1)\T , \qquad f(2,1) = 1 + 1 = 2 \ge 1 \quad\text{— отвергается;}
\]
назад:
\[
  (0,1)\T , \qquad f(0,1) = 1 + 1 = 2 \ge 1 \quad\text{— тоже отвергается,}
\]
координата $x_1$ не меняется.

Координата $x_2$, вперёд:
\[
  (1,2)\T , \qquad f(1,2) = 0 + 0 = 0 < 1 \quad\text{— принимается.}
\]
\[
  x^{(2)} = (1,2)\T , \qquad f\big(x^{(2)}\big) = 0 < 1 .
\]

\medskip
\noindent\textbf{Поиск по образцу.}
\[
  x_{\text{обр}} = 2x^{(2)} - x^{(1)} = (1,3)\T , \qquad f(1,3) = 0 + 1 = 1 > 0 ,
\]
откат, базовая точка $x^{(2)} = (1,2)\T$.

\medskip
Дальше исследующий поиск из $(1,2)\T$ при $h_i = 1$ улучшения не даёт (все
четыре пробные точки дают $f = 1 > 0$), поэтому шаги дробятся:
$h_i := \beta h_i = \tfrac12$, затем $\tfrac14$, $\tfrac18$ и так далее. Каждый
раз поиск безуспешен, пока не выполнится $\max_i h_i < \eps = 0{,}1$, и метод
останавливается в точке
\[
  x_* = (1, 2)\T , \qquad f(x_*) = 0 ,
\]
найденной уже на второй итерации. Значения в базовых точках:
$5 \to 1 \to 0$.
\end{example}

\begin{idea}
Характерная деталь: оба поиска по образцу были отвергнуты. Так получается
потому, что функция симметрична и шаг сетки крупный — удачное направление
$x^{(k)} - x^{(k-1)}$ «перелетает» через минимум. На вытянутых оврагах картина
обратная: именно шаги по образцу дают основное продвижение, а исследующий
поиск только подправляет направление.

Останов здесь определяется не близостью к $x_*$, а измельчением сетки: метод
не знает, что уже стоит в минимуме, и тратит ещё несколько безуспешных
исследующих поисков на дробление $h_i$ до $\eps$.
\end{idea}
